% Также можно использовать \Referat, как в оригинале
\chapter*{Реферат}
Целью данной курсовой работы является разработка алгоритмов эффективной трёхмерной укладки коробок на паллету. Актуальность исследования обусловлена необходимостью максимально компактного размещения грузов для транспортировки, что позволяет снизить логистические затраты и оптимизировать использование складских площадей.

В работе реализованы два алгоритма: жадный GreedySolver (среднее заполнение 73.6\%) и метаэвристический LNSSolver на основе поиска в большой окрестности (Large Neighborhood Search), достигающий заполнения до 90.4\%. Разработана эффективная структура данных HeightHandler для хранения карты высот паллеты — сравнительное тестирование пяти реализаций показало ускорение в 12.9 раза. Реализованы метрики оценки устойчивости укладки и интерактивный 3D-визуализатор.

Тестирование проводилось на 436 реальных задачах от компаний-партнёров. Исходный код доступен в открытом репозитории.

% TODO: Раскомментировать после добавления ссылки на GitHub
% Исходный код: \url{https://github.com/Straple/Palletizer}

Ключевые слова: паллетизация, трёхмерная упаковка, Large Neighborhood Search, складская логистика, NP-трудная задача.

%В работе XX страниц, XX иллюстраций, XX источников.

\chapter*{Abstract}
The purpose of this course work is to develop algorithms for efficient three-dimensional box packing on a pallet. The relevance of this research is determined by the need for the most compact cargo placement for transportation, which reduces logistics costs and optimizes warehouse space utilization.

Two algorithms were implemented: the greedy GreedySolver (average fill rate 73.6\%) and the metaheuristic LNSSolver based on Large Neighborhood Search, achieving up to 90.4\% fill rate. An efficient HeightHandler data structure was developed for storing the pallet height map — comparative testing of five implementations showed a 12.9x speedup. Stability metrics and an interactive 3D visualizer were implemented.

Testing was conducted on 436 real-world tasks from partner companies. The source code is available in an open repository.

% TODO: Uncomment after adding GitHub link
% Source code: \url{https://github.com/Straple/Palletizer}

Keywords: palletization, 3D bin packing, Large Neighborhood Search, warehouse logistics, NP-hard problem.

%The thesis consists of XX pages, XX illustrations, XX references.
\pagestyle{plain} % нумерация вкл.
%%% Local Variables: 
%%% mode: latex
%%% TeX-master: "rpz"
%%% End: 
